% Section 2 of the manuscript: background and related work.
% The cells of the comparison table are filled from the documentation of each tool, with its version and the date
% it was checked; a cell that rests on a published comparison rather than on the tool's own documentation says so.

\section{Background and related work}
\label{sec:related}

\subsection{Keystroke logging in writing research}
\label{sec:related-logging}

Keystroke logging records every key press of a writer with a time, so that the process of writing can be studied
alongside its product \citep{wengelin2023}. From such records researchers compute inter-key intervals and
pauses, often classified by their location in the text (within a word, between words, between sentences), revisions,
bursts of production between pauses, and the amount of text produced and the time taken \citep{conijn2019}. Five programs developed for academic
use and free of charge have been described as available for this purpose: CyWrite, GenoGraphiX-Log, Inputlog,
ScriptLog and Translog-II \citep{wengelin2023}. Inputlog records in any Windows program; the others record in an editor
of their own, which a recent comparison describes as not ecological \citep{crotti2026}. That comparison also asks
whether a logger supports non-alphabetic writing and whether it records key-down and key-up times, and it introduces a
hybrid logger for non-alphabetic scripts \citep{crotti2026}.

Table~\ref{tab:loggers} compares these loggers with the one presented here, adding the dimensions that concern a log
of input-method conversion: how the logger handles the conversion, whether its log holds the text or the names of
keys, and how it was validated. Each cell was filled from the documentation of the tool as published on the date it
was checked; where the documentation examined is silent, the cell says so, and where a cell rests on the published
comparison rather than on the tool's own documentation, it is marked. Translog-II records the key presses that a writer
makes in the input method, together with the characters they produce, in its own editor \citep{zou2026,carl2012}.
Inputlog has test versions for Korean and Chinese that record the input together with the resulting characters
\citep{inputlog-web}. The hybrid logger reconstructs the state of the text by reading it through plug-ins for
particular applications, detects the commits of the input method by inference, and was validated with two participants
and one input method against screen recordings \citep{crotti2026}. Its detection figures are not set against this
study's, because they come from human input and count commits, whereas this study uses synthetic input and matches
intervals. The weaknesses of the present logger are kept in the table: it was not validated with human input or in
applications other than the reference fields, it does not separate inter-key intervals by script, and it records
release times without validating them.

% A landscape long table (main.tex): a sideways float of this length runs off the rotated page.
\begin{landscape}
{\tablebodyfont
\begin{longtable}{L{0.17}L{0.14}L{0.19}L{0.14}L{0.13}L{0.21}}
\caption{Keystroke loggers for writing research, as documented on 5 October 2026}
\label{tab:loggers}\\
\toprule
Logger and version & Where it records & Input-method conversion & Text or key names in the log & Key-down and key-up & Validation reported \\
\midrule
\endfirsthead
\tablecontinued{6}
\toprule
Logger and version & Where it records & Input-method conversion & Text or key names in the log & Key-down and key-up & Validation reported \\
\midrule
\endhead
CyWrite; no tagged release, repository described as a preliminary public release, last updated 21 July 2026 \citep{cywrite-repo} & Its own editor in a web browser \citep{wengelin2023,cywrite-repo} & Not documented & Text changes and keystrokes \citep{wengelin2023} & Yes$^{a}$ & Not examined$^{b}$ \\
GenoGraphiX-Log 2.1.1 (Beta-3), released 15 September 2025 \citep{ggxlog-web} & Its own editor \citep{wengelin2023} & Hardware keyboard and mouse input only; software-driven input not supported \citep{ggxlog-web} & Session log and final text \citep{ggxlog-web} & Yes$^{a}$ & Not documented in the sources examined \\
Inputlog 9.7.0.1, March 2026 \citep{inputlog-web} & Any Windows program \citep{inputlog-web} & Not supported in version 9.5$^{a}$ & Keystrokes \citep{wengelin2023} & Yes in version 9.5$^{a}$ & Not examined$^{b}$ \\
Inputlog Korean and Chinese versions, test releases of March 2025 \citep{inputlog-web} & Windows & Records the input together with the selected or resulting characters \citep{inputlog-web} & Text & Not documented & Not documented in the sources examined \\
ScriptLog; current version platform-independent, available from its developers \citep{wengelin2023}; the website, last updated in 2003, offers version 1.1 for Windows under a paid licence \citep{scriptlog-web} & Its own editor \citep{wengelin2023} & Not supported$^{a}$ & Keystrokes \citep{wengelin2023} & Yes$^{a}$ & Not examined \\
Translog-II 2.24 \citep{translog-web} & Its own editor \citep{wengelin2023} & Records the presses typed into the input method with the characters they produce \citep{zou2026} & Text \citep{zou2026} & Key-down only \citep{zou2026}$^{c}$ & Not documented in the sources examined \\
Hybrid logger; source repository, no release, last updated 13 May 2026 \citep{hylog-repo} & Plug-ins for a word processor and a web browser, with Inputlog \citep{crotti2026} & Commits inferred from keystrokes and text; inter-key intervals before a candidate is confirmed and after a conversion told apart \citep{crotti2026} & Text and keystrokes \citep{crotti2026} & Yes \citep{crotti2026} & Two participants, one input method, against screen recordings \citep{crotti2026} \\
This logger 0.1.0, in release 1.0.0 of the code \citep{release} & Any application, system-wide; the application is not recorded & Windows: display and cloaking of the input method's windows; macOS: inferred from key presses; the time of a commit is approximated by the row's time and duration & Neither (fourteen key categories) & Yes; release times not validated & Synthetic input against reference input fields, three cells, pre-registered; no human input \\
\bottomrule
\end{longtable}
\par\smallskip\noindent\footnotesize
$^{a}$As reported in the comparison by \citet{crotti2026}, which examined Inputlog 9.5, not in the tool's own
documentation.
$^{b}$The articles that report a validation of the tool are not openly accessible and were not read, so the column
reflects only the sources examined.
$^{c}$\citet{crotti2026} describe Translog-II as having no visibility of key-down and key-up events.\par
}
\end{landscape}

\subsection{Writing in Japanese and Chinese}
\label{sec:related-cjk}

Keystroke data have been used to study writing in languages that pass through an input method. A hidden Markov model
of writing processes has been estimated from the keystroke data of high-school and university students in Japan; the
article does not state the language of writing and does not treat the input method \citep{uto2020}. In a thesis on writing Chinese with a keyboard, pauses between typing the Pinyin and selecting the
characters emerged during coding as a category of their own, next to the categories derived from alphabetic writing
\citep{lu2020}. An English-to-Japanese corpus of translation processes recorded with Translog-II is publicly available
\citep{carl2016}, and its logs record the keys typed into the input method together with the text
\citep{tprdb-enja15}. The hybrid logger names Japanese, with its several scripts and input modes, as future work
\citep{crotti2026}. The present study does not analyse writing; it asks whether the conversion of Japanese input can be
logged without the text and how closely such a log agrees with a reference.

\subsection{Keystroke records without text}
\label{sec:related-content}

Keystroke datasets have been released without the typed text, for example with only the hand that pressed each key,
whether it was the space key, and the timings \citep{tappy}. Not recording text is therefore not new in itself; the difference here is a log of the
conversion that holds no text. Keystroke records without text can still reveal it: key press and release timings form a
side channel that yields nontrivial information about typed English words even under substantial measurement error
\citep{monaco2018}, and the temporal features of typing are used to recognise users \citep{whiskerd2020}. This study
estimates how far its key rows reveal the typed phrases (Section~\ref{sec:recoverability}); it judges no record that
keeps times and does not evaluate the identification of a writer.

\subsection{Validation of timing instruments}
\label{sec:related-timing}

Instruments that time responses in behavioural research are validated against an external device. Without external
chronometry a timing error can go unnoticed while the log files suggest that a study ran normally, and benchmarks with
idealised equipment and simplified scripts are unrepresentative of what researchers achieve in the field
\citep{plant2016}. Keystroke sequences have been validated by closing the key switches of a keyboard with such a device
and comparing the recorded intervals \citep{pinet2017}, and a comparison of experiment software warns that its timings
represent something approaching a best case \citep{bridges2020}. The present study places its reference at the
software layer, because its question is agreement with the calls that an application receives from the input method,
which no external device observes; what this gives up is stated in Section~\ref{sec:method-reference}. The timing of
the key rows that this study reports covers the path from the point at which a press is injected into the operating
system to the row, a part of the path that the validations above measured from the key to the record; it is read
together with their figures for the rest, such as the sampling of a USB keyboard every 8~ms, which bounds the precision
of measured intervals \citep{pinet2017}.

\subsection{Methods of agreement}
\label{sec:related-agreement}

Agreement is estimated with limits of agreement \citep{blandaltman1986}, Cohen's kappa \citep{cohen1960}, the
intraclass correlation \citep{shroutfleiss1979,mcgrawwong1996} and percentile bootstrap intervals that resample
whole blocks of trials \citep{efron1979,cameron2008}; their use is described in Section~\ref{sec:method-analysis}.
